% =====================================================================
\section{The wall in two coordinates: what the two pictures jointly prove}
\label{ch:twocoord}
% =====================================================================

\subsection{The picture}

The campaign now has two independent descriptions of the same obstruction. We state the picture
once, plainly, and then say exactly what it does and does not prove.

\textbf{Coordinate 1 (measure / weak-$*$).} The gap is the positive isolated atom $f_1^2>0$ of
the K\"all\'en--Lehmann measure of $\operatorname{tr}F^2$ at the bottom of its support. The atom
is exactly the weak-$*$ discontinuity of the boundary density (Chapter~\ref{ch:weakstar}). Every
finitely computable datum is weak-$*$ continuous, and atomless measures are weak-$*$ dense, so no
forcing functional can be continuous. A proof in this coordinate must be a weak-$*$
discontinuous, acid-clean, $(L,a)$-uniform functional that is provably positive on the
Yang--Mills measure. The sharpest known candidate is the negative-association inequality
$\operatorname{Cov}(G,\Phi'|0)\le0$.

\textbf{Coordinate 2 ($\beta$ / RG / conformal).} The gap is the statement that pure
Yang--Mills ($N_f=0$) lies outside the conformal window, i.e.\ that $\beta_{N_f=0}$ has no
strong-coupling zero (Chapter~\ref{ch:betarg}). The conformal window is real, so
asymptotic freedom alone does not imply a gap; no exact identity forbids an infrared fixed point;
and every modern $\beta$/CFT technique either tests the wrong thing (CFT existence rather than RG
flow) or is phase-blind (central charges) or is compatible with a gapless endpoint (holographic
NEC, DSE scaling). A proof in this coordinate must be a nonperturbative sign theorem for
$\beta_{N_f=0}$ or a transport of the strong-coupling gap across the crossover.

\textbf{The invariant.} In both coordinates the missing object is the same: the failure of the
infrared theory to be a gapless CFT, equivalently the positive isolation of the $0^{++}$ atom,
equivalently $H1$. The paper's structural claim --- a labelled conjecture, not a theorem --- is
that $H1$ and $(G1)$ are the same wall drawn in two coordinates.

\subsection{Why the coordinates are genuinely disjoint}
\label{sec:disjoint}

Disjointness is not rhetoric; it can be stated precisely, and it is the reason the two failures
are independent evidence.

\textbf{The measure argument does not apply to the $\beta$ coordinate.} The weak-$*$ theorem is
about functionals on the space of measures. A beta function is not a measure, and the statement
``$\beta_{N_f=0}$ has no zero'' is not a statement about weak-$*$ continuity. One could try to
encode $\beta$ as a functional of the measure, but that encoding is precisely the dynamical
dictionary whose proof is the gap. So the weak-$*$ obstruction is silent in the $\beta$
coordinate.

\textbf{The window argument does not apply to the measure coordinate.} The existence of the
conformal window is a statement about the massless-fermion content of a gauge theory. It says
nothing about the atom of a given measure. The measure coordinate's obstruction would persist
even if the window were empty; the $\beta$ coordinate's obstruction would persist even if the
measure were atomless by fiat. So the window argument is silent in the measure coordinate.

\textbf{The technique pools are disjoint.} The measure coordinate uses spectral theory, moment
problems, complex analysis, optimal transport, dependence inequalities, and numerical
eigenvalue methods. The $\beta$ coordinate uses the conformal bootstrap, the $a$-theorem,
central charges, large-$N$, holography, Dyson--Schwinger equations, and lattice step-scaling.
The intersection is small (both use lattice simulations, but for different observables), and the
key arguments are of different types: one is an analytic statement about measures, the other a
dynamical statement about flows.

\textbf{The acid tests are different.} In the measure coordinate the acid test is ``does the
functional vanish on the free massless field?'' In the $\beta$ coordinate the acid test is ``does
the argument accidentally also prove a gap for $N_f=12$?'' The first is a statement about
measures; the second is a statement about matter content. A tool that passes one has no reason to
pass the other, and indeed no tool in the campaign passes both.

This is the precise sense in which the two-coordinate picture is more than a restatement: it is
two independent falsification programmes, each of which failed, in different languages, for
different reasons.

\subsection{A third obstruction: the essential singularity}

There is a third obstruction that is neither the weak-$*$ discontinuity nor the dynamical
$\beta$-function problem, although it is related to both: the gap is an essential singularity in
the coupling,
\begin{equation}
m\sim\Lambda\sim\mu\,e^{-1/(2\beta_0g^2)} .
\label{eq:ess3}
\end{equation}
Every derivative vanishes at $g=0$; every analytic-in-$g^2$ method returns zero; Borel
resummation inherits the renormalon ambiguity of the condensate. This is a theorem about the
class of analytic transfers (Chapter~\ref{sec:esssing}). It is not a third coordinate of the same
invariant in the sense of measure-versus-flow; it is a third \emph{no-go class}, cutting across
both coordinates. The campaign records it separately because it explains why the ``solve exactly
then continue'' family fails identically, and because it is the cleanest statement of why the
perturbative expansion cannot manufacture the nonperturbative scale.

The three obstructions are:
\begin{enumerate}
\item \textbf{weak-$*$ discontinuity} (measure coordinate): finite data cannot see the atom;
\item \textbf{essential singularity} (analytic-transfer class): analytic continuation cannot
  manufacture the scale;
\item \textbf{dynamical $\beta$ sign} ($\beta$ coordinate): no exact identity fixes the sign of
  the strong-coupling beta function.
\end{enumerate}
They are not independent in the logical sense --- a proof of the gap would resolve all three ---
but they are independent in the evidentiary sense: each blocks a different technique pool, and
the failure of each is established by different arguments.

\subsection{Independence of the evidence}

The strongest statement the campaign can make is the following, and it is deliberately modest.
\begin{itemize}
\item In the measure coordinate, the obstruction is proved as a theorem about continuous
  functionals (Chapter~\ref{ch:weakstar}, Theorem 4), and the classification of the campaign's
  tools into blind/read/support-only classes is established route by route.
\item In the $\beta$ coordinate, the non-existence of a kinematic proof is proved by the reality
  of the conformal window, and the insufficiency of each modern technique is established by a
  precise argument (Chapter~\ref{ch:betarg}).
\item The two arguments do not use each other. Neither is a corollary of the other. Neither can
  be repaired by importing the other's methods, because the methods do not apply (the
  disjointness of \S\ref{sec:disjoint}).
\end{itemize}
So the wall has been hit twice, by two disjoint toolboxes, with two independent diagnoses. This
is the campaign's strongest evidence that the wall is a property of the problem rather than of a
particular language. It is still not a proof: two failed attacks are not a no-go theorem for all
attacks. But it is far stronger than one failed attack repeated in many costumes, and it explains
why the immune map is so uniform.

\subsection{What a proof needs, stated once}

Combining the two coordinates and the three obstructions, a proof of the gap must supply an
object with the following properties. In measure language:
\begin{enumerate}
\item[(P1)] it is a functional of the field measure, not merely of spectral data (else it is in
  the weak-$*$ continuous class and blind);
\item[(P2)] it is acid-clean: zero on the free massless field and on the atomless gapped
  composite;
\item[(P3)] it is \emph{inhomogeneous} in the coarse mass --- a genuine source, not a
  coefficient --- so that the gapless solution is not a fixed point of the associated
  self-consistency relation;
\item[(P4)] it is nonperturbative and non-circular: valid at $g^2=O(1)$, without assuming the
  gap, threshold, or isolation it is meant to prove;
\item[(P5)] it is $(L,a)$-uniform and survives the continuum limit.
\end{enumerate}
In $\beta$ language, the same object is a nonperturbative proof that $\beta_{N_f=0}$ has no zero
for $g>0$, or equivalently a transport of the strong-coupling gap across the Savvidy crossing
with a controlled continuum limit. The five conditions are the same five conditions of
Chapter~\ref{sec:missing}; the two languages are two ways of writing the same specification.

\subsection{Does two coordinates make it easier?}

No. It makes it sharper. The honest answer has three parts.
\begin{itemize}
\item \textbf{Not easier.} Neither coordinate supplies a new estimate. The measure coordinate
  proves that a large class of tools cannot work; the $\beta$ coordinate proves that a large
  class of techniques cannot decide the question. Neither produces the missing object.
\item \textbf{Sharper.} The two coordinates specify the target in two independent languages, and
  an object satisfying the specification in one language must satisfy it in the other. This is a
  strong consistency condition: a candidate tool can be tested in both coordinates before it is
  believed. The campaign used this as a filter (for example, a candidate that looked like a
  $\beta$-function argument but was secretly a moment inequality failed the measure acid test
  immediately).
\item \textbf{More honest.} The two-coordinate picture explains why the campaign's many failures
  were not failures of effort. They were failures of type: the tools were in the wrong class. The
  picture does not turn a failure into a success; it turns a confusing failure into a diagnosed
  one.
\end{itemize}

\subsection{The one-bit statement}

The slot wave's sentence deserves the last word of this chapter, because it is the campaign's
best summary of the situation. The exact relation that would close the gap is coextensive with
the whole quantum field theory: in the measure coordinate it is the full K\"all\'en--Lehmann
measure, in the $\beta$ coordinate it is the full nonperturbative beta function, and in the
analytic coordinate it is the full transseries. The equation is infinite-dimensional in every
coordinate. But its irreducible core is one scalar sign condition: \emph{is there a sign-definite,
$Q$-independent source at the zero?} As a system, the problem is the theory; as a decision, it is
one bit; and the one bit is the mass gap.

This is the campaign's final structural statement. It is not a proof. It is the precise
description of what a proof would have to be.
